\chapter{Business Requirements Document (BRD)}
\label{ch:brd}

\section{Purpose}

This chapter records the business context that the repository supports. The
repository is a developer tool with no commercial surface, so business
requirements are necessarily modest and are stated only where the code provides
evidence or where an assumption is clearly labelled.

\section{Business problem}

\begin{longtable}{p{3.4cm}p{11cm}}
\caption{Business problems and repository evidence}\\
\toprule
\textbf{Problem} & \textbf{Evidence} \\
\midrule

Repetitive screen construction consumes developer time & Recipe catalog exists precisely to remove repeated screen authoring; \texttt{packages/registry/data/**}. \\
Inconsistent screen quality across a team & A single shared registry plus deterministic rendering produces identical output; \texttt{packages/core/src/generate.ts}. \\
Dependency setup drift & Automated install using the project's package manager; \texttt{packages/cli/src/utils/installDeps.ts}. \\
Screens fall behind catalogue improvements & Versioned recipes, manifest tracking and diff-previewed updates; \texttt{packages/cli/src/commands/update.ts}. \\
\bottomrule
\end{longtable}

\textbf{Status:} INFERRED --- these are the business consequences the design
implies. The repository contains no business case document, no market analysis
and no customer research.

\section{Business objectives}

Each objective below is expressed as something the current code can already do,
so that no objective is aspirational:

\begin{enumerate}
  \item Ship a reusable catalog of best-practice Expo / React Native screens
        (24 recipes present).
  \item Make adding a screen a single, reversible action
        (\texttt{expo-snips add}, with no overwrite of unrelated files).
  \item Adapt generated code to the developer's existing stack rather than
        requiring the developer to adapt (stack detection and conditional
        templates).
  \item Keep installed screens maintainable over time (manifest, \texttt{status},
        \texttt{update}).
  \item Distribute the catalog and packages without manual release steps
        (GitHub Actions workflows).
\end{enumerate}

\textbf{Status:} IMPLEMENTED (objectives 1--5 are all realised in code).

\section{Stakeholders}

\begin{longtable}{p{4.2cm}p{10.2cm}}
\caption{Stakeholders}\\
\toprule
\textbf{Stakeholder} & \textbf{Interest} \\
\midrule

End users (developers) & Adding complete, correct screens quickly and consistently. \\
Recipe authors / contributors & An unambiguous way to define screens (JSON Schema, validation script). \\
Maintainers & Publishing the registry and npm packages reliably; keeping the catalog correct. \\
Organisation / copyright holder & Declared in \texttt{packages/vscode-extension/LICENSE} as ``Jay Jogane'' (2026). No other organisational information exists in the repository. \\
\bottomrule
\end{longtable}

\section{User groups}

\begin{itemize}
  \item \textbf{CLI users} --- interact through \texttt{init}, \texttt{list},
        \texttt{add}, \texttt{status}, \texttt{update}.
  \item \textbf{VS Code users} --- interact through the completion trigger, the
        command palette and the status bar.
  \item \textbf{Catalog consumers} --- anyone fetching \texttt{index.json} or
        individual recipe JSON documents from the registry.
  \item \textbf{Catalog producers} --- maintainers adding recipes and running
        \texttt{pnpm validate} and \texttt{pnpm build:registry}.
\end{itemize}

\textbf{Status:} IMPLEMENTED.

\section{Business processes}

\begin{longtable}{p{4cm}p{10.4cm}}
\caption{Business processes and how they are performed}\\
\toprule
\textbf{Process} & \textbf{Implementation} \\
\midrule

Publish the screen catalog & Push to \texttt{main}/\texttt{master} $\rightarrow$ \texttt{node scripts/build-registry.mjs} $\rightarrow$ deploy the \texttt{registry/} folder to GitHub Pages (\texttt{deploy-registry.yml}). \\
Publish the packages & Push a \texttt{v*} tag $\rightarrow$ \texttt{pnpm install} $\rightarrow$ \texttt{pnpm build} $\rightarrow$ publish \texttt{@expo-snips/core} and \texttt{@expo-snips/cli} to npm (\texttt{publish.yml}). \\
Add a screen recipe & Author a JSON file under \texttt{packages/registry/data/} plus \texttt{.hbs} templates; validate with \texttt{scripts/validate.mjs}; rebuild. \\
Distribute to users & Users install the CLI from npm or the extension from a \texttt{.vsix}; the catalog is bundled at build time and is also fetched from the registry when configured. \\
\bottomrule
\end{longtable}

\section{Business requirements}

\begin{longtable}{p{1.6cm}p{5.4cm}p{7.4cm}}
\caption{Business requirements}\\
\toprule
\textbf{ID} & \textbf{Requirement} & \textbf{Evidence / status} \\
\midrule

BR-01 & The catalog must be distributable independently of the clients. & IMPLEMENTED --- the registry is a folder of static JSON published to GitHub Pages; clients fall back to a bundled copy. \\
BR-02 & Clients must work without network access. & IMPLEMENTED --- bundled \texttt{REGISTRY\_RECIPES} in both the CLI and the extension. \\
BR-03 & Generated code must not overwrite unrelated files. & IMPLEMENTED --- only recipe file paths are written; primitives skip files that already exist (\texttt{primitives.ts}). \\
BR-04 & Installs must be inventoried per project. & IMPLEMENTED --- \texttt{.expo-snips.json} with recipe id, version, options, variant, files and timestamp. \\
BR-05 & Users must be able to review changes before applying updates. & IMPLEMENTED --- \texttt{update} prints a line diff and requires confirmation unless \texttt{--yes}. \\
BR-06 & The catalog must be extensible by third parties. & IMPLEMENTED --- JSON Schema plus a self-hosted registry option (\texttt{--registry}, \texttt{EXPO\_SNIPS\_REGISTRY}, \texttt{expo-snips.registryUrl}). \\
BR-07 & Release of packages must be automated. & IMPLEMENTED --- \texttt{publish.yml} runs on \texttt{v*} tags. \\
BR-08 & The VS Code extension should also be released automatically. & NOT IMPLEMENTED --- the extension is packaged manually with \texttt{vsce} (\texttt{package.json} script ``package''); no workflow builds it. \\
BR-09 & No usage telemetry or analytics. & IMPLEMENTED (by absence) --- no analytics, telemetry or tracking code exists. \\
\bottomrule
\end{longtable}

\section{Expected outcomes}

\begin{itemize}
  \item Developers can add a complete screen with one command.
  \item Generated code matches the project's router, styling, UI primitives, icon
        library and language.
  \item Screens remain updateable from the catalog.
  \item Maintainers can add recipes without changing engine code.
\end{itemize}

\textbf{Status:} IMPLEMENTED.

\section{Success metrics}

\textbf{UNKNOWN.} The repository contains no analytics, telemetry, usage counters,
issue templates, download badges or any other measurable signal. No success
metric can be reported without invention, and none is proposed here as fact.

\textbf{Status:} RECOMMENDED (a future measurement programme would be needed to
establish metrics such as number of generations, recipe coverage or update
adoption --- but this document does not invent targets).

\section{Constraints}

\begin{itemize}
  \item Node.js $\geq$ 18 is required (\texttt{engines.node} in the CLI).
  \item Generated screens consume Expo Router by default; the catalog's
        react-navigation coverage is limited to \texttt{auth/login} and
        \texttt{auth/register}.
  \item The shadcn UI path forces \texttt{styling=nativewind} and
        \texttt{icons=lucide} (\texttt{add.ts}, lines 202--214 and
        \texttt{generator.ts}, lines 82--85).
  \item The registry is unauthenticated static hosting; there is no access
        control on published recipes.
  \item The VS Code extension targets VS Code $\geq$ 1.80.0.
\end{itemize}

\section{Assumptions}

\begin{itemize}
  \item Users have an existing Expo / React Native project and a working package
        manager.
  \item Users will review generated code and replace placeholder business logic
        before shipping.
  \item ``Icons = none'' is an acceptable degradation path; templates guard icon
        imports with \texttt{\{\#if (eq icons "lucide")\}}.
\end{itemize}

\textbf{Status:} INFERRED.

\section{Risks}

\begin{longtable}{p{1.5cm}p{5.2cm}p{1.9cm}p{5.8cm}}
\caption{Business risks}\\
\toprule
\textbf{ID} & \textbf{Risk} & \textbf{Severity} & \textbf{Mitigation / note} \\
\midrule

R-01 & Users treat placeholder auth hooks as production logic. & High & Templates carry a \texttt{// Replace with your actual authentication logic} comment and a commented \texttt{supabase} example; user-facing docs should state this plainly. \\
R-02 & A remote registry is fully trusted: arbitrary recipe templates are rendered to disk and their npm dependencies are executed. & High & No signature or checksum verification exists. Mitigated only by defaulting to the bundled catalog. \\
R-03 & Declared-but-unimplemented options (\texttt{tamagui}, \texttt{vector-icons}, \texttt{custom}) cause user confusion. & Medium & Document explicitly; see Chapter~\ref{ch:known-issues}. \\
R-04 & No automated tests means regressions in the engine or CLI are detectable only manually. & Medium & Add a test suite; see Chapter~\ref{ch:testing}. \\
R-05 & Licensing inconsistency (MIT in two packages, proprietary elsewhere; two different repository URLs). & Medium & Reconcile licence and repository metadata before public distribution. \\
R-06 & A user's existing files at recipe paths are overwritten by \texttt{add}. & Medium & \texttt{add} writes unconditionally via \texttt{fs.writeFile}; only primitives are skip-aware. \\
\bottomrule
\end{longtable}

\textbf{Status:} INFERRED (severity and mitigation are this document's assessment,
grounded in the cited code).
